% Revision 10. Runtime baseline dec1a40f; experiment-specific exceptions
% and the dated review boundary are retained explicitly.
\section{Runtime contract and experimental coverage}
\label{app:runtime}

This appendix describes the runtime at \code{dec1a40f} (8 October 2026),
using the source and run-record review dated 14 October. It does not
assert that every reported experiment ran this revision. Housing and
Procurement used the October executor; Data center used an earlier
executor with a branch-local registry edit; Refund V2.1 used the
September executor. Appendix~\ref{app:runtime-cases} records those paths.
Issue and review statuses below refer to that dated review, not to a
fresh audit of the repository on the date of this revision.

\subsection{Runtime overview}
\label{app:runtime-overview}

Figure~\ref{fig:runtime} shows the execution and receipt path. Families
supply the environment and verifier; the kernel supplies phase
scheduling, evidence records and receipt finalization. Family-specific
drivers still choose entry points, route seats to players and compute
some headline measures. These responsibilities are not yet standardized
(Appendix~\ref{app:runtime-cases}).

\begin{figure}[htbp]
\centering
\begin{tikzpicture}[
  node distance=5mm and 7mm,
  box/.style={draw=ink!70, rounded corners=2pt, align=center,
    font=\small, text width=42mm, minimum height=12mm, inner sep=3pt},
  arr/.style={-{Latex[length=2mm]}, thick, draw=ink!85}]
\node[box] (plan) {World, block, profiles, seeds\\\textbf{resolved plan}};
\node[box,right=of plan] (phase) {Freeze observations\\act, then transition};
\node[box,right=of phase] (log) {Provider/tool events\\and action envelopes};
\node[box,below=of log] (seal) {Reconstruct state, score\\then seal evidence};
\node[box,left=of seal] (receipt) {Receipt with typed scores\\and reference values};
\node[box,left=of receipt] (replay) {Replay state and score\\no model call};
\draw[arr] (plan)--(phase);
\draw[arr] (phase)--(log);
\draw[arr] (log)--(seal);
\draw[arr] (seal)--(receipt);
\draw[arr] (receipt)--(replay);
\end{tikzpicture}
\caption{The receipt path. Replay checks recorded transitions and receipt
scores, not a fresh provider execution, the harness or parsing. It does
not automatically cover a headline analysis computed outside the
receipt. The kernel supplies no sandbox or arbitrary-crash recovery.}
\label{fig:runtime}
\end{figure}

\subsection{Records and identity}
\label{app:runtime-records}

A world is one instance of a case design, represented by a kernel
\code{CaseManifest}. The case studies use \emph{case} for the design as
a whole. Authoring records are immutable dataclasses that reject unknown fields.
Portable case, metric, cluster and cell identifiers follow a lowercase
alphanumeric grammar with internal underscores, periods and hyphens;
colons and whitespace are excluded. This avoids ambiguity with rLLM's
task/rollout separator. A wrapper may normalize an upstream identifier,
retaining the raw identifier in provenance; the kernel itself never
rewrites it.

\paragraph{Canonical bytes.}
Structured records use JSON with sorted keys, compact separators, raw
UTF-8 and no non-finite floats. Artifacts and published files are hashed
as exact bytes. The serializer distinguishes $-0.0$ from $0.0$.
Default-valued optional fields may be omitted from canonical form, so
adding the deferred and inapplicable metric lists did not change older
receipt digests. Digests establish consistency with the bound bytes;
they do not independently attest that a provider's report is true.

\paragraph{Plan resolution.}
A \code{RunSpec} names the suite, evaluation blocks, profiles and seat
assignments. Resolution produces a \code{RunPlan} containing the world
records, sampling plan, blocks, profiles, implementation pins, input
digests and cells, and writes its digest before provider execution.
An implementation pin names a component, kind, version and digest.
A cell includes the world content hash, family and version, block,
world and sampling seeds, cluster identifier and level, optional pair
identifier and seat-to-profile map. A block declares subject seats,
controlled profiles, repetitions, one of controlled, cross-play,
self-play, reference or human comparison, and a paired, independent or
fixed seed policy.

Admission compares declared harness requirements with declared provider
capabilities before spending. At this baseline it also rejects an
OpenRouter profile that combines reasoning effort and reasoning token
budget, a pair the provider refuses. A separate unscored canary probes
live compatibility. Neither check guarantees that a route will remain
available or implement a requested control faithfully.

\paragraph{Profiles and keys.}
An \code{AgentProfile} records the declared model, harness and version,
prompt, tools, memory, runtime settings, reasoning condition, effort and
budget, rationale visibility, sampling, action/time/cost budgets and
retry policy. The reasoning declaration enters the plan and per-seat
profile hashes. It can therefore be checked for an in-place change,
but not for provider compliance. A provider request seed lives in the
profile and is not part of the cell identifier. Identifiers may repeat
across plans; the publication key is the triple (run-plan identifier,
cell identifier, episode-attempt identifier).

\subsection{Phase scheduling}
\label{app:runtime-phases}

A \code{PhaseSpec} names actors, mode (single, sequential or
simultaneous), action cap, successor phases and invalid-action policy.
Under \code{reject}, an invalid action stops the episode with a
scheduler contract error. Under \code{family_defined}, the family
receives an invalid envelope and determines the consequence, such as a
forfeited turn. An exception inside a family hook is a family defect,
not a verdict about the subject.

For a single or simultaneous phase, all observations are computed
before actions are requested, and the transition receives the complete
set of action envelopes. This prevents observation order from revealing
a peer's action within the phase. It does not remove order effects
inside a family's transition. Parsing and legality have typed results;
both an unparsable action and a parseable illegal action retain their
reason in an invalid envelope. The scheduler enforces phase and world
action caps and retains the world seed for seeded transitions and
scoring. The scheduler is provider-free and obtains actions through an
abstract response source.

\subsection{Harnesses, provider calls and retry policy}
\label{app:runtime-harness}

A harness turns observations into actions through an
\code{AttemptContext}, which offers model and tool ports, fixed budget
values, and the attempt identity and seed. The model port seals a
request before sending it and supports native tools, a JSON dialect or
text. An undeclared, ungranted or over-budget tool is stopped before
dispatch and logged as \code{tool_dispatch_rejected}, with reason
\code{undeclared_tool}, \code{tool_not_granted} or
\code{tool_budget_exceeded}. \code{tools_not_admitted} instead means
that a harness received tool calls without any granted tool port.
The budget view supplies the attempt's round allowance, output-token
limit and profile cost cap as fixed values, not continuously updated
remaining spend. Wall-clock timeouts are applied outside the
replayable path and become typed failures.

The default \code{minimal_chat/1.0} harness makes one model call without
tools or memory. Explicit registration is required for
\code{native_tool_chat/1.0}, which supports multiple native tool calls,
and \code{json_dialect/1.0}, which requests a schema-constrained JSON
object. Only the native-tool harness is used for the stated upstream
$\tau^2$-bench comparability claim \citep{barres2025tau2}; the JSON
fallback is labelled non-comparable. Optional LangChain and LangGraph
adapters run behind the same ports for harness comparisons. Each
harness registers open, act, close and failure-classification hooks,
with replayability, blocking and sub-agent requirements.

\paragraph{Request and response records.}
Requests bind the provider, model, instructions, input or native
messages, tool definitions, sampling and reasoning controls, output
limit, timeout, cost limit, output schema and declared seed. A stream
flag enters the digest when set. Only the OpenRouter client streams;
it reassembles chunks into the non-streamed response shape and treats
a stream without a terminal finish reason as a transport failure.
Responses record the reported model, finish reason, empty/truncated
status, input/cached/output tokens and reported cost. Failed calls can
also carry billed tokens and cost. The OpenAI Responses and Arena
clients reject streamed requests at this baseline. Arena is an
OpenAI-compatible client used by Refund, not only an admission probe.

The OpenRouter client checks the reported provider, canonical model and
routing information against the sealed pin, failing on a mismatch. It
cannot verify the report independently or identify an unreported
deployment change. It requires a declared seed; whether the endpoint
honors it is recorded rather than enforced.

\paragraph{Retries and truncation.}
The API clients disable SDK retries. An action condition is retried
only if the profile names it, subject to kernel limits. Configuration,
authentication and account faults (HTTP 402) are never retried.
\code{provider_choice_error} identifies an upstream generation error
and is retryable only by name. An HTTP 404 after a route has succeeded
in the same cell is \code{provider_rejected_after_route_proven}; it is
also retried only when declared. Other route rejections do not acquire
retryability merely because admission succeeded.

A length retry doubles the output budget up to the profile's
\code{max_output_tokens_ceiling}, or eight times the declared budget by
default. No further length retry is issued at the ceiling, and the
widened limit applies to every round of the harness attempt. Empty
completions are rejected. An answer cut off before completion is a
typed \code{length} failure before family parsing; a truncated response
whose answer had already closed is retained. This differs from the
September executor, which could pass an incomplete truncated reply to
the parser after retries were exhausted.

\subsection{The CLI subjects used in the case studies}
\label{app:runtime-cli}

Opus~5.5 runs through the kernel's \code{ClaudeCodePrintClient}; Sol~6.1
runs through \code{CodexExecClient} in \code{tools/cli_subjects}, outside
the kernel package. They are experimental subjects, not merely
diagnostics. In the reviewed panels each action launches one ephemeral
subprocess at low effort. The clients record executable version and
binary digest, not the path, working directory or environment. Claude
Code's binary is re-hashed before each call. Its client defaults to
medium effort if a request omits the control; the reported panels
explicitly set low.

These subjects do not apply seed, temperature or output-token controls.
An output budget present in a profile is therefore a declaration, not
an applied CLI limit. Claude Code runs with tools disabled and no
persistent session. The reviewed calls contain two internal model turns
within one kernel provider call; internal CLI retries are not separate
kernel attempts. Codex retains its tool definitions, and a call that
uses a tool is refused. It reports no served-model identity. Costs are
list-price estimates for subscription logins, not billed charges.

The reviewed executables were Claude Code~2.1.292 and Codex~0.160.0,
with recorded binary digest prefixes \code{97a01e5b} and
\code{112fae7a}; the Codex adapter's digest prefix was \code{3f254873}.
Version and digest pins limit executable drift, but do not certify the
CLI's internal execution or hosted model deployment.

\subsection{Evidence, interruption and publication}
\label{app:runtime-evidence}

Each attempt owns an append-only JSON-lines log and a content-addressed
artifact directory. File locks, exclusive creation and device/inode
checks guard the log. An event records its sequence, type, plan, cell,
episode and attempt identities, phase/action/call identities,
visibility label, payload digest, previous hash and own hash. Calls and
tool invocations have start and terminal events; logical actions also
have an \code{agent_action_failure} terminal. Other events record
dispatch intents and rejections, retry back-off, phases, parse and
legality results, transitions, termination, outcomes and scores. A
bookkeeping failure is an outcome-unknown tool event with the
corresponding condition, not a separate event type.

The start is flushed before the side effect. Finalization reconciles
identities, reconstructs scoring input, runs the scorer, appends the
score and seals event and artifact roots before writing the receipt.
A sealed store rejects later writes. Resume verifies and continues the
chain; it is not arbitrary-crash recovery. An interruption caught by the
running executor records \code{outcome_unknown}. A hard kill may leave
a start without a terminal, which resume does not repair and
reconciliation refuses to finalize. A timeout can also be
outcome-unknown; a declared retry creates a new attempt with the same
request bytes. The log does not prove whether an uncertain external
side effect happened.

\paragraph{Publication and access.}
The hierarchy is campaign, run, task, attempt and event, with the sealed
plan at the run root. Publication selects fields rather than copying
the store wholesale. A manifest binds published file digests and can
record a full \code{source_commit}; a provenance helper locates a
commit containing the pinned source bytes. The commit is manifest
metadata, not an added receipt field.

The baseline publication policy excludes raw responses, full prompts,
model reasoning, complete live receipts and failure messages. Live
bundles carry sanitized receipt projections with source-receipt, world
and plan digests; provider-free golden receipts may be published in
full. A named-field token scan checks the declared publication policy,
but is not a semantic secret detector. Public bundles alone therefore
do not expose all evidence needed for replay. A reader with the
unpublished run root can perform the audit in
Appendix~\ref{app:runtime-replay}; a public reader can check the published
files and their bindings, but cannot reconstruct omitted evidence from
a digest.

\subsection{Receipts and the measurement contract}
\label{app:runtime-receipts}

Each declared metric is a \emph{leaf}. Its \code{EvaluationReceipt}
records a measurement status, either \code{ok} or
\code{invalid_measurement}, and a separate inclusion status. It also
binds plan/world/cell/block/episode/attempt identities,
sampling and pairing information, declared per-seat profile hashes,
implementation pins, evidence seal, primary metric, score envelopes,
optional failure, observability limits and replay level. At most one
receipt is included per planned cell. This cardinality check cannot
prove that the attempt was selected before seeing its outcome; that is
a campaign procedure.

A \code{ScoreEnvelope} identifies the metric, verifier kind, reference
kind, input scope and direction (maximize, minimize or none), and
contains its value, further metrics, reference values, validity report
and evidence references. Multi-seat scores may include utility and
capture by seat. The five verifier names in Section~\ref{sec:measurement} map to
\emph{answer key} (\code{canonical_reference}), a declared correct
resolution; \emph{rule} (\code{rule_constraint}), an outcome or path
constraint; \emph{value} (\code{objective_reference}), economic value
against a policy or bound; \emph{relative} (\code{comparative}), standing
against named agents or populations; and \emph{rubric}
(\code{rater_judge}), an artifact rated against a rubric. The paper additionally states
each check's output and outcome/path scope; these labels do not imply
that every legacy family stores the same fields in the same way.
Outputs are pass/fail, counts or amounts. Checks remain in their native
units; rubric results are reported beside computed checks.

Value references distinguish a reachable policy, the best policy given
the decider's information, and a full-information upper bound. A
feasible policy is a witness, not a guaranteed outcome floor. An
information-privileged planner is not the attainable policy of a seat
that lacks that information. Relative checks must name the counterpart
population. Answer keys, rules and rubric raters require their own
validity checks; recording a typed envelope does not supply them.

\paragraph{Declared metric policy.}
The measurement manifest assigns each leaf a finalize-time or deferred
scope, a cell or subject-seat scope, an explicit reduction where
multiple subject seats are combined, and case-conditional applicability.
Deferral names the awaited artifact before the run. At finalization,
applicable finalize-time metrics must be returned, deferred metrics
must be listed, and inapplicable conditional metrics must be listed
separately. Undeclared, omitted applicable or returned inapplicable
metrics are violations. Deferred and case-conditional metrics cannot
be the primary or an admission metric. These checks run at
finalization, replay and audit.

The common scoring-contract test exercises enrolled scorers on witness
fixtures, including a mutation that changes each trajectory-derived
metric. Bridge-gated fixtures run only when the pinned upstream bridge
is provisioned; continuous integration provisions three of eight such
families. A passing local suite with skipped bridge assertions is not
upstream validation. Registration, fixture enrollment and actual
receipt production are reported separately below.

\subsection{Replay and reproducibility}
\label{app:runtime-replay}

Replay invokes the family's phase, initial-state, actor, transition,
termination, outcome and scoring hooks. For each action it reads the
recorded observation, request, canonical response, parse result and
legality result. It compares state boundaries and scores, stopping on
divergence. It does not rerun the harness, parser or provider. The
caller supplies the sealed plan and registry through an evaluation
setup; subject seats and profiles are read from that plan.

The finalizer assigns \code{state_and_score} to every completed scored
execution it seals, including one excluded as
\code{invalid_measurement}. Operational-failure receipts have replay
level \code{none}. The schema reserves \code{score_only}, which the
current finalizer does not emit. A replayed exclusion is not a newly
scored successful episode.

\begin{samepage}
At the October baseline, the reviewer command
\begin{quote}\small
\begin{verbatim}
aeread verify-replay BUNDLE --run-root R
  --setup MODULE:CALLABLE
\end{verbatim}
\end{quote}
\end{samepage}
rechecks published receipt bindings using the retained run evidence.
It reports verified, differing or missing-evidence rows and exits
unsuccessfully unless its checks pass. Coverage is
\code{every_published_episode} only when the bundle includes the
trajectory grain identifying those episodes; otherwise it is limited
to receipts found under the run root. Operational-failure receipts are
verified as exclusions. This command replaces reliance on a campaign's
self-reported replay boolean; its existence does not show that it ran
on every panel. The Housing v14 reports are local, Procurement's panels
were checked inline, and the Data center branch predates the command.

Seeds, request digests and evidence roots support reconstruction of a
recorded execution, not deterministic regeneration by a model.
Provider serving can vary across batches and deployments
\citep{he2025nondeterminism}. The OpenRouter routing check binds the
provider's report to the declared route, not to an independently
attested deployment. The five \code{execution_mode} values are
validated metadata; the invoked entry point and retained evidence,
not that field alone, determine whether execution is live or replayed.

\subsection{Family plugins, adapters and enrollment}
\label{app:runtime-families}

A family registers an exact identifier and version. Duplicate
registration and missing required hooks are rejected. Thirteen hooks
cover validation, initial state, phases, actors, observation, parsing,
legality, transition, termination, outcome, scorer/reference
construction and generation; optional hooks cover teardown and
case-conditional metrics. The registry has no entry-point discovery or
dynamic import, and the kernel has no static family-package dependency.
The reviewer command is an explicit exception: its
\code{--setup MODULE:CALLABLE} imports the caller-named family setup.

The inventory has ten wrappers of upstream code, two paper-based
re-implementations and nine native/fixture packages. Eight wrappers
execute upstream code in fidelity tests. AucArena instead tests
hand-copied functions; STEER's answer-key suite is adapter-owned
because the pinned upstream removed its scorer. AgenticPay runs the
upstream bridge and replays goldens but lacks a standalone parity
recomputation. Algorithmic collusion and TERMS-Bench are explicitly
labelled re-implementations. A parity check measures wrapper fidelity,
not model performance or construct validity.

Fifteen packages have the common scoring-contract fixture. Five
built-ins do not: retail, the two data-center packages, consent and the
smoke fixture. The data-center packages nevertheless finalize and
publish receipts. Refund V2.1 runs through contributed registration,
calls the same executor and finalizer, and publishes kernel receipts;
it is outside the trusted list and lacks the common fixture. At the
September snapshot the reviewed Refund bundles already contained
1,320 receipt rows across seven bundles, and data-center bundles held
476 rows, 405 fully replayed. The 19 legacy summaries carrying both
replay flags were data-center-terms campaigns, not evidence of their
exclusion from the receipt path.

For retail, the static parity check compares the initial database,
ordered tool calls and results, final database and diff, database reward,
communication check, judge inputs and recorded verdict, and total
reward against upstream gold trajectories: 114 tasks and an 18-task
pilot. That check is not yet rerun by CI and is not a live-model result.
Three bridge fidelity jobs (AgenticPay, AmazonHistoryPrice and GovSim)
run in CI; a flag makes a skipped required bridge test fail.

Existing evaluation tools provide trajectory logging and multi-agent
interfaces \citep{ukaisi2024inspect,terry2021pettingzoo,
lanctot2019openspiel,agapiou2022meltingpot}. We describe AERead's
integration of seat assignments, economic references and replay
obligations without claiming those individual capabilities are absent
elsewhere. Harbor export is planned, not implemented. The experimental
rLLM integration wraps the separate exchange-economy engine, not this
kernel \citep{laude2026harbor,agentica2025rllm}.

\begingroup
\small
\setlength{\tabcolsep}{4pt}
\renewcommand{\arraystretch}{1.13}
\begin{longtable}{@{}P{1.44in}P{1.77in}P{2.64in}P{0.55in}@{}}
\caption{Family-package inventory at \code{dec1a40f}. Check names match
Section~\ref{sec:measurement}; the primary is listed first. Static
world counts describe the package snapshot, not the later panels in
Table~\ref{tab:case-paths}.}\label{tab:families}\\
\toprule
\textbf{Family} & \textbf{Origin and seats} & \textbf{Checks and instances} & \textbf{Test}\\
\midrule\endfirsthead
\multicolumn{4}{@{}l}{\textit{Table \thetable\ continued: family inventory}}\\
\toprule
\textbf{Family} & \textbf{Origin and seats} & \textbf{Checks and instances} & \textbf{Test}\\
\midrule\endhead
\midrule\multicolumn{4}{r}{\textit{Continued on the next page}}\\\endfoot
\bottomrule\endlastfoot
\multicolumn{4}{@{}l}{\textit{Wrapped upstream benchmarks}}\\[2pt]
AgenticPay & Wrapped; 2 seats & Value, rule; 28 static worlds & $\checkmark$\\
Alympics WAC & Wrapped; 5 seats & Relative, rule; 7 worlds & $\checkmark$\\
AmazonHistoryPrice & Wrapped; 2 seats & Relative, rule, value; 45 worlds. Relative scalar primary. & $\checkmark$\\
AucArena & Wrapped components; 1 model / 2 rule bidders & Relative, rule; 6 worlds & $\checkmark$\\
EconAgent & Wrapped; $N=4$ or $10$ seats & Rule (accounting identity; no direction), relative; 3 scenarios & $\checkmark$\\
EconEvals & Wrapped; 1 seat with tools & Value, rule; 28 worlds & $\checkmark$\\
GovSim & Wrapped; 5 seats & Relative, rule; 9 worlds & $\checkmark$\\
NegotiationArena & Wrapped; 2 seats & Relative, rule; 6 worlds & $\checkmark$\\
STEER & Corpus import; 1 seat & Answer key; 1,595 instances & $\checkmark$\\
$\tau$-bench retail & Wrapped; 1 seat / simulated user & Answer key; deferred rubric secondary; 114 tasks & $\circ$\\
\addlinespace[5pt]
\multicolumn{4}{@{}l}{\textit{Re-implemented from a paper}}\\[2pt]
Algorithmic collusion & Re-implemented; 2 seats & Relative, rule, answer key; 6 worlds & $\checkmark$\\
TERMS-Bench & Re-implemented; 1 seat / scripted counterpart & Rule, relative; 30 worlds & $\checkmark$\\
\addlinespace[5pt]
\multicolumn{4}{@{}l}{\textit{Native families and fixture}}\\[2pt]
Housing & Native; campaign default 6 tenants / 4 landlords & Value (upper bound); generated worlds & $\checkmark$\\
Procurement allocation & Native; 1 seat & Value (upper bound); 149 static worlds & $\checkmark$\\
Procurement grounding & Native; 1 seat & Answer key; 1 world & $\checkmark$\\
Commercial-state calibration & Native; 1 seat & Answer key; 9 worlds & $\checkmark$\\
Data-center development & Native; 1 model / 1--5 scripted seats; 5 ids, 10 versions & Value / answer key; 117 worlds. Publishes receipts. & $\circ$\\
Data-center development terms & Native; 1 seat & Answer key; 22 worlds. Publishes receipts. & $\circ$\\
Consent / individual rationality & Native; 1 seat & Value (exact); 1 world & $\circ$\\
Refund & Native; V1.3: 2 seats; V2.1: 3 model-eligible / 1 scripted & Value (V1.3), rule (V2.1); 18 static worlds; V2.1 generated. Produces kernel receipts. & C\\
Single offer & Smoke fixture; 1 seat & No primary verifier; 0 static worlds & $\circ$\\
\end{longtable}
\endgroup

{\small\textbf{Legend.} $\checkmark$: trusted built-in with the common
scoring fixture (15 packages). $\circ$: trusted built-in without that
fixture (five). C: contributed registration (Refund). Symbols describe
enrollment, not receipt availability: Data center and Refund publish
receipts. Retail has a finalizer path but no published live receipts in
this snapshot. The separate exchange engine is excluded; package counts
are distinct from family ids.\par}

\begingroup
\small
\setlength{\tabcolsep}{4pt}
\renewcommand{\arraystretch}{1.13}
\begin{longtable}{@{}P{1.34in}P{2.22in}P{2.96in}@{}}
\caption{Pinned provenance and fidelity evidence for the inventory.
Upstream names are repository identifiers; hashes are commit prefixes.
Fidelity evidence concerns the adapter, not an evaluated model.}
\label{tab:family-fidelity}\\
\toprule
\textbf{Family} & \textbf{Source and license} & \textbf{Fidelity evidence}\\
\midrule\endfirsthead
\multicolumn{3}{@{}l}{\textit{Table \thetable\ continued: provenance and fidelity}}\\
\toprule
\textbf{Family} & \textbf{Source and license} & \textbf{Fidelity evidence}\\
\midrule\endhead
\midrule\multicolumn{3}{r}{\textit{Continued on the next page}}\\\endfoot
\bottomrule\endlastfoot
AgenticPay \citep{liu2026agenticpay} & SafeRL-Lab/AgenticPay\newline\code{1ff4e1a2}; MIT & Byte-identical replay of five goldens through upstream; CI bridge job. No standalone parity recomputation.\\
Alympics WAC \citep{mao2023alympics} & microsoft/Alympics\newline\code{caed7c8c}; MIT & Adapter compared with a second upstream instance.\\
AmazonHistoryPrice \citep{xia2024bargaining} & TianXiaSJTU/\allowbreak AmazonPriceHistory\newline\code{834ad906}; Apache-2.0 & Byte-identical golden replay; CI bridge job.\\
AucArena \citep{chen2023aucarena} & jiangjiechen/\allowbreak auction-arena\newline\code{d0f3bc85}; Apache-2.0 & Component checks and mutation tests over vendored hand-copied functions; upstream LLM bidder not executed.\\
EconAgent \citep{li2024econagent} & tsinghua-fib-lab/\allowbreak ACL24-EconAgent\newline\code{bfada09}; no root license; embedded files BSD-3-Clause & Per-agent state parity across all three scenarios.\\
EconEvals \citep{fish2025econevals} & sara-fish/\allowbreak econ-evals-paper\newline\code{e1f2a40f}; MIT & Byte-identical canonical replay.\\
GovSim \citep{piatti2024govsim} & giorgiopiatti/govsim\newline\code{1d11adf0}; MIT & Adapter compared with raw upstream; CI bridge job.\\
NegotiationArena \citep{bianchi2024negotiationarena} & vinid/NegotiationArena\newline\code{c447fafd}; MIT & Byte-identical upstream-direct replay.\\
STEER \citep{raman2024steer} & narunraman/STEER\newline\code{d66673c8}; no license file & Adapter-owned deterministic answer-key suite; pinned upstream had removed its scorer.\\
$\tau$-bench retail \citep{barres2025tau2} & sierra-research/tau2-bench\newline\code{fc0055dc} (v1.0.1); MIT & Static parity receipt over 114/114 tasks and 18/18 pilot gold trajectories; not rerun in CI.\\
Algorithmic collusion \citep{fish2024collusion} & arXiv:2404.00806v6; HTML and PDF digests pinned; no code & Closed-form calculations compared with the paper appendix.\\
TERMS-Bench \citep{zhang2026termsbench} & arXiv:2605.13909v2; paper digest pinned; upstream link unavailable & Replay on two scripted episodes.\\
Housing & Native & Closed-form assignment oracle for the inventory's family baseline; the v14 reference is described in the case studies.\\
Procurement allocation & Native & Optimal script attains the bound for the inventory baseline; the pay-or-test information-constrained policy is separately identified.\\
Procurement grounding & 231-project snapshot & No additional fidelity evidence listed.\\
Commercial-state calibration & Sanitized archive & No additional fidelity evidence listed.\\
Data-center development & Native & No additional inventory-level fidelity evidence listed; agreement panel replay is in Table~\ref{tab:case-paths}.\\
Data-center development terms & Pinned SEC EDGAR filings & Provenance-pinning test.\\
Consent / individual rationality & Native & No additional fidelity evidence listed.\\
Refund & Native & Internal tests; contributed kernel path. V2.5 is a later provider-free release, outside this inventory snapshot.\\
Single offer & Native smoke fixture & No fidelity claim.\\
\end{longtable}
\endgroup


\subsection{Execution paths behind the case-study results}
\label{app:runtime-cases}

Table~\ref{tab:case-paths} separates the executor from case-specific
drivers and headline analysis. Hashes are abbreviated digest prefixes,
not interchangeable with source commit identifiers. The October
\code{execution.py} prefix is \code{607a79ce}, shared by
\code{6aaf46bb} and \code{dec1a40f}; the September prefix is
\code{3cd3278c}. A later baseline description does not retrospectively
upgrade an older run's behavior or evidence.

\begingroup
\small
\setlength{\tabcolsep}{4pt}
\renewcommand{\arraystretch}{1.16}
\begin{longtable}{@{}P{1.28in}P{2.18in}P{3.06in}@{}}
\caption{What produced the case-study numbers, as documented by the
14 October review. Inline replay, a later independent audit and
headline-measure reconstruction are separate checks.}
\label{tab:case-paths}\\
\toprule
\textbf{Panel} & \textbf{Executor, driver and client} & \textbf{Receipt audit and headline coverage}\\
\midrule\endfirsthead
\multicolumn{3}{@{}l}{\textit{Table \thetable\ continued: experimental execution paths}}\\
\toprule
\textbf{Panel} & \textbf{Executor, driver and client} & \textbf{Receipt audit and headline coverage}\\
\midrule\endhead
\midrule\multicolumn{3}{r}{\textit{Continued on the next page}}\\\endfoot
\bottomrule\endlastfoot
Housing v14\newline 1,800 planned cells; separate 48-cell fill
& October executor \code{607a79ce}; run at \code{9441abca}, 6 October. Family price driver; OpenRouter on pinned routes, through the seat router; minimal-chat harness.
& Original panel: 1,764 completed receipts and 36 operational-failure receipts. Inline replay, then bundle replay of 600 rows per model (564 scored for Gemini), and 48 fill rows; reports retained locally. The original certificate scores six tenants, not the headline seat-0 measure. A later sealed-event derivation is separately versioned.\\
\addlinespace[4pt]
Housing CLI\newline 20 worlds, 40 cells per subject
& October executor \code{607a79ce}; \code{housing_cli.py} on the Housing branch. Claude Code and the external Codex client use the same seat router.
& 40 completed cells per subject; inline receipt replay only. No published bundle or independent bundle-audit report in the review. Inherits Housing's seat and headline-score gaps. The driver is not digested in the plan.\\
\addlinespace[4pt]
Procurement CLI\newline 24 worlds, two briefs, two subjects
& October executor \code{607a79ce}; \code{procurement_cli.py}; runs on 9 October after \code{e60878f8} and \code{f1ca35ec}. Same CLI clients as Housing.
& 24 receipts per run root; inline equality replay. No published bundle or independent bundle-audit report. Headline decision loss is computed from an unsealed action projection by an unpinned reference module. Retained run roots are outside the repository.\\
\addlinespace[4pt]
Procurement flash\newline 24 worlds per model
& October kernel; \code{pay_or_test_openrouter.py}, digest \code{f5ec607c}, committed at \code{d91fbde8}. Runs 13 October; driver bytes match the later commit. Pinned OpenRouter clients.
& 24 completed receipts per model; 72 and 76 provider calls. Inline replay. An offline verify mode exists, but the review found no output from it. Same headline decision-loss gap as the CLI panels.\\
\addlinespace[4pt]
Data-center agreement\newline 18 bundles; 2,456 negotiations
& September executor \code{3cd3278c}; branch based on \code{7ac08f0e}; runs 9--12 October. \code{agreement_cli.py} freezes 11 sources. Branch-local registry edit \code{03a7faef} was not in those frozen sources. CLI subjects and zero-cost scripted clients.
& 2,417 valid negotiations; inline receipt-equality replay. Joint value lost, seat surpluses and IR violations are receipt leaves. Published replay booleans come from the running driver; this branch predates independent bundle replay. The registry edit lacked non-author review.\\
\addlinespace[4pt]
Refund V2.1\newline four models, 120 cells each in the reported panel
& September kernel, identical to \code{9aa11a2d}; \code{v2_runner}; Arena client. Contributed registration calls the shared executor and finalizer.
& Published projections carry \code{state_and_score}; inline replay. Gap figures are regenerated by \code{v2_gap --check}. Produces kernel receipts despite being outside the trusted scoring-contract fixture.\\
\addlinespace[4pt]
Refund V2.5\newline 816 provider-free trajectories
& \code{24c7edd0}, based on \code{dec1a40f}, without kernel changes. Contributed QC path; all seats use a scripted provider; paid calls refused.
& Release replays and audits all 816 receipts. No published bundle in the review. These runs qualify release conformance, not paid-model performance; they do not replace the V2.1 results.\\
\end{longtable}
\endgroup


\paragraph{Housing attribution and headline scoring.}
Housing v14's plan assigns the subject profile to all six tenant seats,
but a driver-side \code{SeatRouterClient} supplies scripts for seats
1--5. Request and receipt profile hashes therefore name a model for
seats it never played; the actual routing map is in an unsealed side
log (HL-T-09). The original receipt's primary score is the six-tenant
sum, not the headline deciding-tenant measure (HL-D-09). For example,
world 300000, pooled landlord, GLM has a receipt sum of
$-1{,}405.86$ and a seat-0 value of $+235.24$. The original headline
analysis reads outcome facts copied to a separate unsealed result file.
The later \code{price_reply_posterior} measure (\code{765e99a7}) derives
the corrected measure from sealed events; that derived analysis must
be identified separately from the original receipt certificate.

\paragraph{Procurement's reference calculation.}
The receipt scores contribution margin and regret against a
full-information bound. The headline decision loss against the best
policy given the buyer's information instead replays an
\code{action_trace} copied to an unsealed result row, through
\code{pay_or_test_reference}, which the original plans do not pin
(P-T-17). This projection drops message fields. The numerical
calculation can be reproduced from those inputs, but the receipt audit
does not certify it. The cross-case replay claim was corrected in
RES-J-04. Standardizing the case-to-runtime interface,
including the subject-seat map and headline metric, is the proposal in
issue~\#289; it is not an implemented guarantee of these runs.

\subsection{Campaign procedure, review and implementation scope}
\label{app:runtime-qc}

The runtime enforces the declared record and measurement contract for
the path it executes. It does not prove that every headline number is
a receipt metric, that a driver routed the declared player, or that a
selected attempt was chosen before its outcome. \qcref{} reports the
checks actually applied: Gate~0, profile; Gate~1, task distribution;
Gate~2, environment and verifier; Gate~3, construct and baselines;
Gate~4, attribution and controls; Gate~5, confirmatory reliability and
publication. The four current designs do not fully pass Gates~1--5;
several profiles and reference-run checks were written after execution.
This is a qualification boundary, not something that a receipt hash
resolves.

A frozen control changed in place must receive a new campaign identity.
The incident log's 8 September 2026 judgment entry, ``a frozen control
edited in place under a published campaign'' (PR~\#98), documents a
digest mismatch after the primary-outcome list changed. The sealed
bundle remained unchanged and the old declaration was restored through
an explicit parameter. Separately, P-J-02 records parent bundles
re-sealed in place and restored after CI detected their changed pinned
digests. These are specific checks and incidents, not evidence that
every campaign satisfied the full procedure.

Eight kernel-lane PRs (\#232--\#239, tracked by \#226) merged without a
non-author approval under a one-time waiver (RN-O-01). The post-merge
review \#266 was still open in the dated review. Passing tests do not
substitute for that missing review. At \code{dec1a40f}, the source count
is 20,747 kernel lines in 31 Python modules, 155,481 family-package
lines, 279 test files and 3,156 test functions. Parametrized collection
counts are omitted because they depend on the checkout and optional
dependencies. These are implementation-size counts, not measures of
benchmark validity or fresh test results for this revision.

\subsection{Limitations and mechanisms motivated by incidents}
\label{app:runtime-limits}

\begin{itemize}
\item \textbf{No sandbox or complete mediation.} Families and tools run
  in-process or through bridge interpreters. The kernel supplies no
  container and does not prevent code from touching resources outside
  its ports. Its logging guarantees cover the instrumented path, not
  all possible external activity.
\item \textbf{Declarations do not verify reasoning.} The OpenRouter and
  Claude Code clients do not populate the typed
  \code{reasoning_tokens} field; such counts occur only in retained raw
  responses. In the first 100 cells per Housing v14 model, the reviewed
  counts averaged 287 for Gemini at minimal, 60 for GLM at low, and 186
  for Luna at low (HL-J-12). Effort labels do not order deliberation
  across models. An earlier route returned about 13 reasoning tokens
  under several declared limits versus 259 with the block omitted;
  that historical incident does not characterize every October route.
\item \textbf{Empty and truncated answers.} Empty completions once
  passed as no-ops and silently dropped about 40\% of a candidate's
  turns, leading to a retracted effect. The port now rejects them.
  October's incomplete-answer handling differs from September's;
  issue~\#265 tracks an erratum for eight earlier procurement rows
  (TB-D-03). Current behavior does not repair those rows retrospectively.
\item \textbf{Run halting is opt-in.} The kernel offers
  \code{max_consecutive_operational_failures}: an adopting sequential
  driver stops after the declared streak or immediately on an account
  fault and marks unreached cells \code{not_attempted}. No family
  driver had adopted that helper at \code{dec1a40f} (\#264). Housing
  used driver-level rules; twelve workers each reaching three failures
  account for its 36 refused cells. There is no kernel campaign-wide
  cost circuit breaker.
\item \textbf{Spend is qualified, not uniformly a lower bound.}
  Provider-reported cost or profile pricing on reported tokens costs a
  call. The October executor charges a billed call even when it later
  fails. The spend reader labels an event-log total \code{exact} or
  \code{lower_bound}, the latter when some cost is unknown; family
  drivers had not adopted that reader at the baseline. CLI list-price
  estimates are a separate quantity. Earlier totals could omit billed
  failed calls.
\item \textbf{Receipt integrity depends on its checker.} An earlier
  compatibility reader coerced falsy values to empty lists, and a writer
  accepted pretty-printed bytes its reader rejected. Fixes and
  regression tests landed before the September snapshot. The incident
  motivates testing the verifier itself; a digest alone is not a
  correctness argument.
\item \textbf{Driver and selection gaps remain.} Headline-metric and
  seat-routing differences are documented above. A critical-failure
  classifier based on exception-message substrings can change stopping
  behavior when wording changes. The pre-outcome attempt-selection
  rule is recorded but not proven by the one-receipt cardinality check.
\item \textbf{Missingness can depend on the subject.} In ten earlier
  Housing campaigns, 47 of 48 trajectory failures involved one model
  family across four providers. That historical concentration limits
  a missing-at-random interpretation; it is not the denominator of
  Housing v14's separate account-fault incident.
\item \textbf{Replay access and recovery are partial.} Independent
  bundle replay requires retained run roots, and some panels have only
  inline checks. A hard-killed attempt may not reconcile. Neither
  arbitrary-crash continuation nor fresh-model determinism is claimed.
\item \textbf{Coverage is version-specific.} The September inventory's
  live evidence does not describe all October panels, which used
  several OpenRouter routes and two CLI subjects. Wrapper parity,
  fixture enrollment, receipts and live-model evidence remain separate
  claims. The original exchange-economy engine runs outside this kernel;
  Refund does not.
\end{itemize}
